%=========================================
% PROOF OF THEOREM 3
%=========================================

\section{Proof of Theorem 3}
\label{Proof of th:feasible}

\subsection{Parts (i)}

Given definition \ref{def: Ztilde}, and because of lemma \ref{lem:trimerror}:
\begin{align*}
&\widehat{Z}_{T,N,R} \\
&=\frac{\mathcal{N}}{\sqrt{T-R}}\sum_{t=R+1}^{T}\frac{1}{N_{\tau(t)}}\sum_{i=1}^{N_{\tau(t)}} r_{i,t} \mathbbm{1} \left\{ \widehat{S}_{i,\tau(t),R}\leq \widehat{S}_{\tau(t),R}^{p(1),L} \right\} \left( \mathbbm{1} \left\{ S_{i,\tau(t)} \leq \widehat{S}_{\tau(t),R}^{p(1)} \right\} +  \mathbbm{1} \left\{ \widehat{S}_{\tau(t),R}^{p(1)} \leq S_{i,\tau(t)} \right\} \right) \\
&- \frac{\mathcal{N}}{\sqrt{T-R}}\sum_{t=R+1}^{T}\frac{1}{N_{\tau(t)}}\sum_{i=1}^{N_{\tau(t)}} r_{i,t} \mathbbm{1}\left\{  \widehat{S}_{\tau(t),R}^{p(\mathcal{N}-1),U} \leq \widehat{S}_{i,\tau(t),R}\right\} \left( \mathbbm{1} \left\{ S_{i,\tau(t)} \leq  \widehat{S}_{\tau(t),R}^{p(\mathcal{N}-1)} \right\} + \mathbbm{1} \left\{ \widehat{S}_{\tau(t),R}^{p(\mathcal{N}-1)} \leq S_{i,\tau(t)} \right\} \right) \\
&=\frac{\mathcal{N}}{\sqrt{T-R}}\sum_{t=R+1}^{T}\frac{1}{N_{\tau(t)}}\sum_{i=1}^{N_{\tau(t)}} r_{i,t} \mathbbm{1} \left\{ \widehat{S}_{i,\tau(t),R}\leq \widehat{S}_{\tau(t),R}^{p(1),L} \right\} \mathbbm{1} \left\{ S_{i,\tau(t)} \leq \widehat{S}_{\tau(t),R}^{p(1)} \right\} \\
&- \frac{\mathcal{N}}{\sqrt{T-R}}\sum_{t=R+1}^{T}\frac{1}{N_{\tau(t)}}\sum_{i=1}^{N_{\tau(t)}} r_{i,t} \mathbbm{1}\left\{ \widehat{S}_{\tau(t),R}^{p(\mathcal{N}-1),U} \leq \widehat{S}_{i,\tau(t),R}\right\} \mathbbm{1} \left\{ \widehat{S}_{\tau(t),R}^{p(\mathcal{N}-1)} \leq S_{i,\tau(t)} \right\} + o_{a.s.}(1)
\end{align*}
Adding and subtracting:
\begin{align*}
&\widehat{Z}_{T,N,R} \\
&=\frac{\mathcal{N}}{\sqrt{T-R}}\sum_{t=R+1}^{T}\frac{1}{N_{\tau(t)}}\sum_{i=1}^{N_{\tau(t)}}r_{i,t}\left( \mathbbm{1}\left\{ \widehat{S}_{i,\tau(t),R}\leq \widehat{S}_{\tau(t),R}^{p(1),L}\right\} \mathbbm{1}\left\{ S_{i,\tau(t)}\leq S_{\tau(t)}^{p(1)}\right\} \right. \\
&\hspace{5.2cm} \left. -\mathbbm{1}\left\{ \widehat{S}_{\tau(t),R}^{p(\mathcal{N}-1),U} \leq \widehat{S}_{i,\tau(t),R} \right\} \mathbbm{1}\left\{ S_{\tau(t)}^{p(\mathcal{N}-1)} \leq S_{i,\tau(t)} \right\} \right) \\
&+\frac{\mathcal{N}}{\sqrt{T-R}}\sum_{t=R+1}^{T}\frac{1}{N_{\tau(t)}}\sum_{i=1}^{N_{\tau(t)}}r_{i,t} \mathbbm{1}\left\{ \widehat{S}_{i,\tau(t),R}\leq \widehat{S}_{\tau(t),R}^{p(1),L}\right\} \\
&\hspace{5.cm} \left( \mathbbm{1}\left\{ S_{i,\tau(t)}\leq \widehat{S}_{\tau(t),R}^{p(1)}\right\} -\mathbbm{1}\left\{ S_{i,\tau(t)}\leq S_{\tau(t)}^{p(1)}\right\} \right) \\
&-\frac{\mathcal{N}}{\sqrt{T-R}}\sum_{t=R+1}^{T}\frac{1}{N_{\tau(t)}}\sum_{i=1}^{N_{\tau(t)}}r_{i,t} \mathbbm{1}\left\{ \widehat{S}_{\tau(t),R}^{p(\mathcal{N}-1),U} \leq \widehat{S}_{i,\tau(t),R} \right\} \\
&\hspace{5.2cm} \left( \mathbbm{1}\left\{ \widehat{S}_{\tau(t),R}^{p(\mathcal{N}-1)} \leq S_{i,\tau(t)} \right\} -\mathbbm{1}\left\{ S_{\tau(t)}^{p(\mathcal{N}-1)} \leq S_{i,\tau(t)} \right\} \right) \\
&= I_{T,N,R}+II_{T,N,R}+III_{T,N,R}.
\end{align*} \hl{Checked until here}

{\small We begin to study $II_{T,N,R}.$ Letting }$\widehat{u}_{\lfloor
t/\tau \rfloor ,R}^{N_{\tau(t)}/\mathcal{N}}=\widehat{S}%
_{\tau(t),R}^{N_{\tau(t)}/\mathcal{N}%
}-S_{\tau(t)}^{p(1)}$%
\begin{eqnarray*}
&&II_{T,N,R}=\frac{\mathcal{N}}{\sqrt{T-R}}\sum_{t=R+1}^{T}\frac{1}{%
N_{\tau(t)}}\sum_{i=1}^{n_{\lfloor t/\tau \rfloor
}}r_{i,t}\left( 1\left\{ \widehat{S}_{i,\tau(t),R}\leq 
\widehat{S}_{\tau(t),R}^{L,N_{\tau(t)}/%
\mathcal{N}}\right\} \right.  \\
&&\phantom{\frac{\mathcal{N}}{\sqrt{T-R}}\sum_{t=R+1}^{T}\frac{1}{n_{\lfloor
t/\tau \rfloor }}\sum_{i=1}^{N_{\tau(t)}}r_{i,t}\left(
\right.}\left. \left( 1\left\{ S_{i,\tau(t)}\leq \widehat{S}%
_{\tau(t),R}^{N_{\tau(t)}/\mathcal{N}%
}\right\} -1\left\{ S_{i,\tau(t)}\leq S_{\lfloor t/\tau
\rfloor }^{N_{\tau(t)}/\mathcal{N}}\right\} \right) \right) 
\\
&=&\frac{\mathcal{N}}{\sqrt{T-R}}\sum_{t=R+1}^{T}\frac{1}{n_{\lfloor t/\tau
\rfloor }}\sum_{i=1}^{N_{\tau(t)}}r_{i,t}\left( 1\left\{ 
\widehat{S}_{i,\tau(t),R}\leq \widehat{S}_{\lfloor t/\tau
\rfloor ,R}^{L,N_{\tau(t)}/\mathcal{N}}\right\} \right.  \\
&&\phantom{\frac{\mathcal{N}}{\sqrt{T-R}}\sum_{t=R+1}^{T}\frac{1}{n_{\lfloor
t/\tau \rfloor }}\sum_{i=1}^{N_{\tau(t)}}r_{i,t}\left(
\right.}\left. \left( 1\left\{ S_{i,\tau(t)}\leq S_{\lfloor
t/\tau \rfloor }^{N_{\tau(t)}/\mathcal{N}}+\widehat{u}%
_{\tau(t),R}^{N_{\tau(t)}/\mathcal{N}%
}\right\} -1\left\{ S_{i,\tau(t)}\leq S_{\lfloor t/\tau
\rfloor }^{N_{\tau(t)}/\mathcal{N}}\right\} \right) \right) 
\\
&=&\frac{\mathcal{N}}{\sqrt{T-R}}\sum_{t=R+1}^{T}\frac{1}{n_{\lfloor t/\tau
\rfloor }}\sum_{i=1}^{N_{\tau(t)}}r_{i,t}\left( 1\left\{ 
\widehat{S}_{i,\tau(t),R}\leq \widehat{S}_{\lfloor t/\tau
\rfloor ,R}^{L,N_{\tau(t)}/\mathcal{N}}\right\} \right.  \\
&&\phantom{\frac{\mathcal{N}}{\sqrt{T-R}}\sum_{t=R+1}^{T}\frac{1}{n_{\lfloor
t/\tau \rfloor }}\sum_{i=1}^{N_{\tau(t)}}r_{i,t}\left(
\right.}\left( \left( 1\left\{ S_{i,\tau(t)}\leq S_{\lfloor
t/\tau \rfloor }^{N_{\tau(t)}/\mathcal{N}}+\widehat{u}%
_{\tau(t),R}^{N_{\tau(t)}/\mathcal{N}%
}\right\} -F\left( S_{\tau(t)}^{N_{\tau(t)}/%
\mathcal{N}}+\widehat{u}_{\tau(t),R}^{n_{\lfloor t/\tau
\rfloor }/\mathcal{N}}\right) \right) \right.  \\
&&\left. -\left( 1\left\{ S_{i,\tau(t)}\leq S_{\lfloor
t/\tau \rfloor }^{N_{\tau(t)}/\mathcal{N}}\right\} -F\left(
S_{\tau(t)}^{p(1)}\right)
\right) \right)  \\
&&+\frac{\mathcal{N}}{\sqrt{T-R}}\sum_{t=R+1}^{T}\frac{1}{n_{\lfloor t/\tau
\rfloor }}\sum_{i=1}^{N_{\tau(t)}}r_{i,t}\left( 1\left\{ 
\widehat{S}_{i,\tau(t),R}\leq \widehat{S}_{\lfloor t/\tau
\rfloor ,R}^{L,N_{\tau(t)}/\mathcal{N}}\right\} \right.  \\
&&\left. \left( F\left( S_{\tau(t)}^{n_{\lfloor t/\tau
\rfloor }/\mathcal{N}}+\widehat{u}_{\tau(t),R}^{n_{\lfloor
t/\tau \rfloor }/\mathcal{N}}\right) -F\left( S_{\lfloor t/\tau \rfloor
}^{N_{\tau(t)}/\mathcal{N}}\right) \right) \right)  \\
&=&\mu _{1}\frac{\mathcal{N}}{\sqrt{T-R}}\sum_{t=R+1}^{T}\left( \left(
1\left\{ S_{i,\tau(t)}\leq S_{\lfloor t/\tau \rfloor
}^{N_{\tau(t)}/\mathcal{N}}+\widehat{u}_{\lfloor t/\tau
\rfloor ,R}^{N_{\tau(t)}/\mathcal{N}}\right\} -F\left(
S_{\tau(t)}^{p(1)}+%
\widehat{u}_{\tau(t),R}^{N_{\tau(t)}/%
\mathcal{N}}\right) \right) \right.  \\
&&\left. -\left( 1\left\{ S_{i,\tau(t)}\leq S_{\lfloor
t/\tau \rfloor }^{N_{\tau(t)}/\mathcal{N}}\right\} -F\left(
S_{\tau(t)}^{p(1)}\right)
\right) \right)  \\
&&+\mu _{1}\frac{\mathcal{N}}{\sqrt{T-R}}\sum_{t=R+1}^{T}\left( F\left(
S_{\tau(t)}^{p(1)}+%
\widehat{u}_{\tau(t),R}^{N_{\tau(t)}/%
\mathcal{N}}\right) -F\left( S_{\tau(t)}^{n_{\lfloor t/\tau
\rfloor }/\mathcal{N}}\right) \right) +o_{p}(1) \\
&=&\mu _{1}\frac{\mathcal{N}}{\sqrt{T-R}}\sum_{t=R+1}^{T}\left( \widehat{S}%
_{\tau(t),R}^{N_{\tau(t)}/\mathcal{N}%
}-S_{\tau(t)}^{N_{\tau(t)}/\mathcal{N}%
}\right) +o_{p}(1),
\end{eqnarray*}%
where the last inequality follows from stochastic equicontinuity, and $\mu
_{1}=\mathrm{E}\left( r_{i,t}1\left\{ S_{i,\tau(t)}\leq
S_{\tau(t)}^{q_{1}}\right\} \right) $ is defined as in the
statement of the Theorem. Since, $III_{T,N,R}$ can be treated by the same
argument as $II_{T,N,R},$ it follows that%
\begin{eqnarray}
&&\widehat{Z}_{T,N,R}  \notag \\
&=&\frac{\mathcal{N}}{\sqrt{T-R}}\sum_{t=R+1}^{T}\frac{1}{n_{\lfloor t/\tau
\rfloor }}\sum_{i=1}^{N_{\tau(t)}}r_{i,t}\left( 1\left\{ 
\widehat{S}_{i,\tau(t),R}\leq \widehat{S}_{\lfloor t/\tau
\rfloor ,R}^{L,(\mathcal{N}-1)N_{\tau(t)}/\mathcal{N}%
}\right\} 1\left\{ S_{i,\tau(t)}\leq S_{\lfloor t/\tau
\rfloor }^{N_{\tau(t)}/\mathcal{N}}\right\} \right.   \notag
\\
&&\phantom{+\frac{\mathcal{N}}{\sqrt{T-R}}\sum_{t=R+1}^{T}\frac{1}{n_{%
\tau(t)}}\sum_{i=1}^{n_{\lfloor t/\tau \rfloor
}}r_{i,t}\left( \right.}\left. -1\left\{ \widehat{S}_{i,\lfloor t/\tau
\rfloor ,R}\geq \widehat{S}_{\tau(t),R}^{U,(\mathcal{N}%
-1)N_{\tau(t)}/\mathcal{N}}\right\} 1\left\{ S_{i,\lfloor
t/\tau \rfloor }\geq S_{\tau(t)}^{(\mathcal{N}-1)n_{\lfloor
t/\tau \rfloor }/\mathcal{N}}\right\} \right)   \notag \\
&&+\mu _{1}\frac{\mathcal{N}}{\sqrt{T-R}}\sum_{t=R+1}^{T}\left( \widehat{S}%
_{\tau(t),R}^{N_{\tau(t)}/\mathcal{N}%
}-S_{\tau(t)}^{N_{\tau(t)}/\mathcal{N}%
}\right) -\mu _{\mathcal{N}}\frac{\mathcal{N}}{\sqrt{T-R}}%
\sum_{t=R+1}^{T}\left( \widehat{S}_{\tau(t),R}^{(\mathcal{N}%
-1)N_{\tau(t)}/\mathcal{N}}-S_{\tau(t)}^{(%
\mathcal{N}-1)N_{\tau(t)}/\mathcal{N}}\right)   \notag \\
&=&Z_{T,N}+\mu _{1}\frac{\mathcal{N}}{\sqrt{T-R}}\sum_{t=R+1}^{T}\left( 
\widehat{S}_{\tau(t),R}^{N_{\tau(t)}/%
\mathcal{N}}-S_{\tau(t)}^{N_{\tau(t)}/%
\mathcal{N}}\right)   \notag \\
&&-\mu _{\mathcal{N}}\frac{\mathcal{N}}{\sqrt{T-R}}\sum_{t=R+1}^{T}\left( 
\widehat{S}_{\tau(t),R}^{(\mathcal{N}-1)n_{\lfloor t/\tau
\rfloor }/\mathcal{N}}-S_{\tau(t)}^{(\mathcal{N}%
-1)N_{\tau(t)}/\mathcal{N}}\right)   \notag \\
&&-\frac{\mathcal{N}}{\sqrt{T-R}}\sum_{t=R+1}^{T}\frac{1}{n_{\lfloor t/\tau
\rfloor }}\sum_{i=1}^{N_{\tau(t)}}r_{i,t}\left( 1\left\{ 
\widehat{S}_{\tau(t),R}^{L,(\mathcal{N}-1)n_{\lfloor t/\tau
\rfloor }/\mathcal{N}}\leq \widehat{S}_{i,\tau(t),R}\leq 
\widehat{S}_{\tau(t),R}^{(\mathcal{N}-1)n_{\lfloor t/\tau
\rfloor }/\mathcal{N}}\right\} 1\left\{ S_{i,\tau(t)}\leq
S_{\tau(t)}^{N_{\tau(t)}/\mathcal{N}%
}\right\} \right.   \notag \\
&&-\left. 1\left\{ \widehat{S}_{\tau(t),R}^{(\mathcal{N}%
-1)N_{\tau(t)}/\mathcal{N}}\leq \widehat{S}_{i,\lfloor
t/\tau \rfloor ,R}\leq \widehat{S}_{\tau(t),R}^{U,(\mathcal{N%
}-1)N_{\tau(t)}/\mathcal{N}}\right\} 1\left\{ S_{i,\lfloor
t/\tau \rfloor }\leq S_{\tau(t)}^{n_{\lfloor t/\tau \rfloor
}/\mathcal{N}}\right\} \right)   \label{Z}
\end{eqnarray}
Recall that, $\widehat{S}_{\tau(t),R}^{(\mathcal{N}%
-1)N_{\tau(t)}/\mathcal{N}}-\widehat{S}_{\lfloor t/\tau
\rfloor ,R}^{L,(\mathcal{N}-1)N_{\tau(t)}/\mathcal{N}}=c_{%
\mathcal{N},\tau(t),R}^{U},$ where $c_{\mathcal{N},\lfloor
t/\tau \rfloor ,R}^{U}=\sup_{i\text{.s.t. }I_{i,T,R}^{(\mathcal{N}%
)}=1}c_{i,\tau(t),R}^{U}.$ We now show that the last term on
the RHS of the last equality in (\ref{Z}) is $o_{p}(1).$ The term reads as%
\begin{eqnarray}
&&\frac{\mathcal{N}}{\sqrt{T-R}}\sum_{t=R+1}^{T}\frac{1}{n_{\lfloor t/\tau
\rfloor }}\sum_{i=1}^{N_{\tau(t)}}\left( r_{i,t}1\left\{ 
\widehat{S}_{\tau(t),R}^{N_{\tau(t)}/%
\mathcal{N}}-c_{1,\tau(t),R}^{U}\leq \widehat{S}_{i,\lfloor
t/\tau \rfloor ,R}\leq \widehat{S}_{\tau(t),R}^{n_{\lfloor
t/\tau \rfloor }/\mathcal{N}}\right\} \right.   \notag \\
&&\phantom{+\frac{\mathcal{N}}{\sqrt{T-R}}\sum_{t=R+1}^{T}\frac{1}{n_{%
\tau(t)}}\sum_{i=1}^{n_{\lfloor t/\tau \rfloor
}}r_{i,t}\left( \right.}\left. -r_{i,t}1\left\{ \widehat{S}_{\lfloor t/\tau
\rfloor ,R}^{(\mathcal{N}-1)N_{\tau(t)}/\mathcal{N}-}\leq 
\widehat{S}_{i,\tau(t),R}\leq \widehat{S}_{\lfloor t/\tau
\rfloor ,R}^{(\mathcal{N}-1)N_{\tau(t)}/\mathcal{N}}+c_{%
\mathcal{N},\tau(t),R}^{U}\right\} \right) ,  \label{BTNR}
\end{eqnarray}%
Let%
\begin{equation*}
C_{1,\tau(t),R}=\left\{ i:1\left\{ \widehat{S}_{\lfloor
t/\tau \rfloor ,R}^{N_{\tau(t)}/\mathcal{N}}-c_{1,\lfloor
t/\tau \rfloor ,R}^{U}\leq \widehat{S}_{i,\tau(t),R}\leq 
\widehat{S}_{\tau(t),R}^{N_{\tau(t)}/%
\mathcal{N}}\right\} =1\right\} ,
\end{equation*}%
that is $C_{1,\tau(t),R}$ is the set of all stocks for which 
$\widehat{S}_{\tau(t),R}^{N_{\tau(t)}/%
\mathcal{N}}-c_{1,\left\lfloor t/\tau \right\rfloor ,R}^{U}\leq \widehat{S}%
_{i,\tau(t),R}\leq \widehat{S}_{\lfloor t/\tau \rfloor
,R}^{N_{\tau(t)}/\mathcal{N}}.$

\noindent Given, Assumption (\ref{ass:crosssec}), the cardinality of $%
C_{1,\tau(t),R}$ is $O_{a.s.}\left( n_{\lfloor t/\tau
\rfloor ,R}^{1+\varepsilon }c_{\mathcal{N},\lfloor t/\tau \rfloor
,R}^{U}\right) .$ Define $C_{\mathcal{N}-1,\tau(t),R}$
analogously. Thus,%
\begin{eqnarray}
&&+\frac{\mathcal{N}}{\sqrt{T-R}}\sum_{t=R+1}^{T}\frac{1}{n_{\lfloor t/\tau
\rfloor }}\sum_{i=1}^{N_{\tau(t)}}\left( r_{i,t}1\left\{ 
\widehat{S}_{\tau(t),R}^{N_{\tau(t)}/%
\mathcal{N}}-c_{1,\tau(t),R}^{U}\leq \widehat{S}_{i,\lfloor
t/\tau \rfloor ,R}\leq \widehat{S}_{\tau(t),R}^{n_{\lfloor
t/\tau \rfloor }/\mathcal{N}}\right\} \right.   \notag \\
&&\left. -r_{i,t}1\left\{ \widehat{S}_{\tau(t),R}^{(\mathcal{%
N}-1)N_{\tau(t)}/(\mathcal{N}-1)}\leq \widehat{S}_{i,\lfloor
t/\tau \rfloor ,R}\leq \widehat{S}_{\tau(t),R}^{(\mathcal{N}%
-1)N_{\tau(t)}/\mathcal{N}}+c_{\mathcal{N},\lfloor t/\tau
\rfloor ,R}^{U}\right\} \right)   \notag \\
&=&\frac{\mathcal{N}}{\sqrt{T-R}}\sum_{t=R+1}^{T}\frac{1}{n_{\lfloor t/\tau
\rfloor }}\sum_{i\in C_{1,\tau(t),R}}^{{}}\left( r_{i,t}-\mu
_{i}\right) -\frac{\mathcal{N}}{\sqrt{T-R}}\sum_{t=R+1}^{T}\frac{1}{%
N_{\tau(t)}}\sum_{i\in C_{\mathcal{N}-1,\lfloor t/\tau
\rfloor ,R}}^{{}}\left( r_{i,t}-\mu _{i}\right)   \notag \\
&&+\frac{\mathcal{N}}{\sqrt{T-R}}\sum_{t=R+1}^{T}\frac{1}{n_{\lfloor t/\tau
\rfloor }}\sum_{i\in C_{1,\tau(t),R}}^{{}}\mu _{i}-\frac{%
\mathcal{N}}{\sqrt{T-R}}\sum_{t=R+1}^{T}\frac{1}{N_{\tau(t)}}%
\sum_{i\in C_{\mathcal{N}-1,\tau(t),R}}^{{}}\mu _{i}  \notag
\\
&=&O_{p}\left( \frac{n_{\tau(t),R}^{1+\varepsilon }\left(
c_{1,\tau(t),R}^{U}+c_{\mathcal{N},\lfloor t/\tau \rfloor
,R}^{U}\right) }{\inf_{t}N_{\tau(t)}}\right)   \notag \\
&&+\frac{\mathcal{N}}{\sqrt{T-R}}\sum_{t=R+1}^{T}\frac{1}{n_{\lfloor t/\tau
\rfloor }}\sum_{i\in C_{1,\tau(t),R}}^{{}}\mu _{i}-\frac{%
\mathcal{N}}{\sqrt{T-R}}\sum_{t=R+1}^{T}\frac{1}{N_{\tau(t)}}%
\sum_{i\in C_{\mathcal{N}-1,\tau(t),R}}^{{}}\mu _{i},
\label{BTNR2}
\end{eqnarray}%
and the last term on the RHS of (\ref{BTNR2}) is zero under $H_{0}$ and of
order $\sqrt{T-R}\frac{n_{\tau(t),R}^{1+\varepsilon }\left(
c_{1,\tau(t),R}^{U}+c_{\mathcal{N},\lfloor t/\tau \rfloor
,R}^{U}\right) }{\inf_{t}N_{\tau(t)}}$ under $H_{A}.$ Them 
{\small both the statement in (i) and in (ii) then follow.}

\subsection{Part (ii)}

Under $H_{A},$%
\begin{eqnarray*}
\widehat{Z}_{T,N,R} &=&Z_{T,N,R}+m_{1}\left( \widehat{S}_{\lfloor t/\tau
\rfloor ,R}^{N_{\tau(t)}/\mathcal{N}}-S_{\lfloor t/\tau
\rfloor }^{N_{\tau(t)}/\mathcal{N}}\right) -m_{\mathcal{N}%
}\left( \widehat{S}_{\tau(t),R}^{(\mathcal{N-}1)n_{\lfloor
t/\tau \rfloor }/\mathcal{N}}-S_{\tau(t)}^{(\mathcal{N-}%
1)N_{\tau(t)}/\mathcal{N}}\right)  \\
&&+O_{a.s.}\left( \sqrt{T-R}\frac{n_{\lfloor t/\tau \rfloor
,R}^{1+\varepsilon }\left( c_{1,\tau(t),R}^{U}+c_{\mathcal{N}%
,\tau(t),R}^{U}\right) }{\inf_{t}N_{\tau(t)}}%
\right) .
\end{eqnarray*}%
Since $Z_{T,N,R}$ diverges at rate $\sqrt{T-R},$ $\widehat{Z}_{T,N,R}$ will
diverge at least at rate $\sqrt{T-R}.$

\section{Proof of Corollary 1}
\label{Proof of cor:error}

Let $q(1)=1/\mathcal{N}$,
\begin{align*}
&\frac{\mathcal{N}}{\sqrt{T-R}} \sum_{t=R+1}^{T} m_{1}\left( \widehat{S}_{\tau(t),R}^{p(1)} - S_{\tau(t)}^{p(1)}\right) \\
&=\frac{\mathcal{N}}{\sqrt{T-R}} \sum_{t=R+1}^{T} m_{1}\left( \widehat{S}_{\tau(t),R}^{\left \lfloor p(1) \right \rfloor }- S_{\tau(t)}^{\left\lfloor p(1)\right\rfloor }\right) +o_{p}(1) \conditionaltext{\textcolor{red}{\text{Understand step}}} \\
&=\frac{\mathcal{N}}{\sqrt{T-R}} \sum_{t=R+1}^{T} m_{1} \left( \widehat{S}_{\tau(t),R}^{q(1)} - \widetilde{S}_{\tau(t)}^{q(1)} \right) +o_{p}(1), 
\end{align*}
where $\widehat{S}_{\tau(t),R}^{q(1)}$ and $\widetilde{S}_{\tau(t)}^{q(1)}$ are the unconditional estimator of the $q(1)$ quantile based on the estimated and the true sorting variable, % See definition of quantile regression (check function)
\begin{align*}
\widehat{S}_{\tau(t),R}^{q(1)} = \arg \min_{S} \frac{1}{N_{\tau(t)}} \sum_{i=1}^{N_{\tau(t)}} \left( q(1) - \mathbbm{1} \left\{ \widehat{S}_{i,\tau(t),R} \leq S \right\} \right) \left( \widehat{S}_{i,\tau(t),R}- S \right), % \label{haS1} 
\end{align*}
\begin{align*}
\widetilde{S}_{\tau(t)}^{q(1)}=\arg \min_{S}\frac{1}{N_{\tau(t)}} \sum_{i=1}^{N_{\tau(t)}} \left( q(1) - \mathbbm{1} \left\{ S_{i,\tau(t)} \leq S \right\} \right) \left( S_{i,\tau(t)}-S\right). 
\end{align*}
Note that the $o_{p}(1)$ on the RHS of the last line in the first equation comes from the fact that, given assumption \crefdefpart{ass:exogenous, i},\footnote{Notice that the unconditional estimator of a quantile is based on a continuum of observations, while an order statistic is based on a finite number $N$ of observations. So the difference between a continuous measure and a discrete one based on $N$ observations would go to the quantile at the speed $1/N$.}
\begin{align*}
\sup_{t\geq R}\left\vert \widehat{S}_{\tau(t),R}^{\left\lfloor p(1)\right\rfloor }-\widehat{S}_{\tau(t),R}^{q(1)}\right\vert =O_{p}\left( \frac{1}{N}\right),
\end{align*}
\begin{align*}
\sup_{t\geq R}\left\vert S_{\tau(t)}^{\left\lfloor p(1)\right\rfloor }-\widetilde{S}_{\tau(t)}^{q(1)}\right\vert =O_{p}\left( \frac{1}{N} \right).
\end{align*}

Let $S^{q(1)}$ be the true unconditional $q(1)$ quantile of $S_{i,\tau(t)}$, which does not depend on time, given assumption \crefdefpart{ass:mixing, i}. By Taylor expanding the pseudo score of $\widehat{S}_{\tau(t),R}^{q(1)}$ around $S^{q(1)}$, up to a term of smaller probability order, 
\begin{align*}
\widehat{S}_{\tau(t),R}^{q(1)}-S^{q(1)} &= \frac{1}{N_{\tau(t)}}\sum_{i=1}^{N_{\tau(t)}}\left( q(1)-\mathbbm{1} \left\{ \widehat{S}_{i,\tau(t),R}\leq S^{q(1)}\right\} \right) \conditionaltext{\textcolor{red}{\text{Minus missing? Taylor expand derivative}}} \\
&= \frac{1}{N_{\tau(t)}}\sum_{i=1}^{N_{\tau(t)}}\left( q(1)-\mathbbm{1} \left\{ S_{i,\tau(t)}\leq S^{q(1)}\right\} \right)  \\
& +\frac{1}{N_{\tau(t)}}\sum_{i=1}^{N_{\tau(t)}}\left( \mathbbm{1} \left\{ \widehat{S}_{i,\tau(t),R}\leq S^{q(1)}\right\} - \mathbbm{1} \left\{ S_{i,\tau(t)}\leq S^{q(1)}\right\} \right) \conditionaltext{\textcolor{red}{\text{or signs inverted}}} \\
&= \frac{1}{N_{\tau(t)}}\sum_{i=1}^{N_{\tau(t)}} \left( q(1)-\mathbbm{1} \left\{ S_{i,\tau(t)}\leq S^{q(1)}\right\} \right) \\
& +\frac{1}{N_{\tau(t)}}\sum_{i=1}^{N_{\tau(t)}} \Big( \Big( \mathbbm{1} \left\{ S_{i,\tau(t)}\leq S^{q(1)}-\left( \widehat{S}_{i,\tau(t),R}-S_{i,\tau(t) }\right) \right\} \\
& \hspace{2cm} -F \left( S^{q(1)}-\left( \widehat{S}_{i,\tau(t),R}-S_{i,\tau(t)}\right) \right) \Big) \\
& \hspace{2cm} - \Big( \mathbbm{1} \left\{ S_{i,\tau(t)}\leq S^{q(1)}\right\} -\underbrace{F\left( S^{q(1)}\right) }_{q(1)} \Big) \Big)  \\
& +\frac{1}{N_{\tau(t)}} \sum_{i=1}^{N_{\tau(t)}} \Big( F\left( S^{q(1)}-\left( \widehat{S}_{i,\tau(t),R}-S_{i,\tau(t)}\right) \right) -\underbrace{F\left( S^{q(1)}\right) }_{q(1)} \Big)  \\
&= I_{T,N} + II_{T,N} + III_{T,N}
\end{align*}
Notice that 
\begin{align*}
I_{T,N}=\left( \widetilde{S}_{\tau(t)}^{q(1)}-S^{q(1)}\right) \left( 1+o_{p}(1)\right) \conditionaltext{\textcolor{green}{Why?}}
\end{align*}
and that 
\begin{align*}
III_{T,N}=\frac{1}{N_{\tau(t)}}\sum_{i=1}^{N_{\tau(t)}}f\left(S^{q(1)}\right) \left( \widehat{S}_{i,\tau(t),R}-S_{i,\tau(t)}\right) +o_{p}(1) \conditionaltext{\textcolor{green}{\text{Sign inverted}}}
\end{align*}
Moreover, given assumption \ref{ass:sups}, because of stochastic equicontinuity, \conditionaltext{\textcolor{green}{Why assumption 5?}}
\begin{align*}
\frac{\mathcal{N}}{\sqrt{T-R}} \sum_{t=R+1}^{T} II_{T,N} &= \frac{\mathcal{N}}{\sqrt{T-R}} \sum_{t=R+1}^{T} \frac{1}{N_{\tau(t)}} \sum_{i=1}^{N_{\tau(t)}} \left( \mathbbm{1} \left\{ S_{i,\tau(t)}\leq S^{q(1)}-\left( \widehat{S}_{i,\tau(t),R}-S_{i,\tau(t) }\right) \right\} \right. \\
& \hspace{5cm} \left. -F\left( S^{q(1)}-\left( \widehat{S}_{i,\tau(t),R}-S_{i,\tau(t)}\right) \right) \right) \\
& + \frac{\mathcal{N}}{\sqrt{T-R}} \sum_{t=R+1}^{T} \frac{1}{N_{\tau(t)}} \sum_{i=1}^{N_{\tau(t)}} \Big( F\left( S^{q(1)}\right) - \mathbbm{1} \left\{ S_{i,\tau(t)}\leq S^{q(1)}\right\} \Big) \\
&= o_{p}\left( I_{T,N} + III_{T,N} \right). \conditionaltext{\textcolor{green}{o_p(1)+o_p(1); Two pieces missing}}
\end{align*}
Thus,
\begin{align*}
&\frac{\mathcal{N}}{\sqrt{T-R}}\sum_{t=R+1}^{T}m_{1}\left( \widehat{S}_{\tau(t),R}^{q(1)}-\widetilde{S}_{\tau(t)}^{q(1)}\right)  \\
&= \frac{\mathcal{N}}{\sqrt{T-R}}\sum_{t=R+1}^{T}m_{1}\left( \widehat{S}_{\tau(t),R}^{q(1)}-S^{q(1)}\right) -\frac{\mathcal{N}}{\sqrt{T-R}}\sum_{t=R+1}^{T}m_{1}\left( \widetilde{S}_{\tau(t)}^{q(1)}-S^{q(1)}\right) \\
&= \frac{\mathcal{N}}{\sqrt{T-R}}\sum_{t=R+1}^{T}\frac{1}{N_{\tau(t)}}\sum_{i=1}^{N_{\tau(t)}}m_{1} f\left(S^{q(1)}\right) \left( \widehat{S}_{i,\tau(t),R}-S_{i,\tau(t)}\right) +o_{p}(1), \textcolor{green}{\text{Understand}}
\end{align*}
and so from theorem \ref{th:feasible} it follows that:
\begin{align*}
\widehat{Z}_{N,T,R}-Z_{N,T} &= \frac{\mathcal{N}}{\sqrt{T-R}}\sum_{t=R+1}^{T} \frac{1}{N_{\tau(t)}}\sum_{i=1}^{n_{\tau(t)}} m_{1}f\left( S^{q(1)}\right) \left( \widehat{S}_{i,\tau(t),R}-S_{i,\tau(t)}\right) \\
& - \frac{\mathcal{N}}{\sqrt{T-R}}\sum_{t=R+1}^{T} \frac{1}{N_{\tau(t)}}\sum_{i=1}^{n_{\tau(t)}} m_{\mathcal{N}}f\left( S^{q(\mathcal{N})}\right) \left( \widehat{S}_{i,\tau(t),R}-S_{i,\tau(t)}\right) 
%\label{SER}
\end{align*}
If assumption \crefdefpart{ass:mixing, v} and \crefdefpart{ass:sups, iii}\conditionaltext{\textcolor{green}{needed?}}and/or $(T-R)/R\rightarrow 0$, the term on the RHS of the last equation is $o_{p}(1)$, \conditionaltext{\textcolor{green}{Why?}}and the statement in (\ref{No-PEE}) holds. If instead assumption \ref{ass:strongdep} holds and $(T-R)/R\rightarrow \pi$, where $0 < \pi < \infty$, then the term on the RHS of the last equation converges in distribution to a normal distribution with zero mean,\conditionaltext{\textcolor{green}{Why?}}and the statement in (\ref{PEE}) follows. 
